\chapter{Methodology}
\label{ch:methodology}

This chapter demonstrates equations and tables -- the two things methodology
chapters use most.

\section{Equations}
\label{sec:equations}

% A numbered, labelled display equation. Refer to it with \ref.
The governing relationship is given in Equation~\ref{eq:example}:
\begin{equation}
    E = mc^2 .
    \label{eq:example}
\end{equation}

% An UNnumbered display equation uses \[ ... \].
The mean of $n$ samples is
\[
    \bar{x} = \frac{1}{n} \sum_{i=1}^{n} x_i .
\]

% Several aligned lines use the align environment from amsmath. The &
% marks the alignment point; \\ ends each line. Each line can have its
% own \label.
\begin{align}
    a &= b + c, \label{eq:line1}\\
    d &= e - f. \label{eq:line2}
\end{align}

Equations~\ref{eq:line1} and~\ref{eq:line2} can be referenced individually.

\section{Experimental Setup}
\label{sec:setup}

% A basic table. Column types: l (left), c (centre), r (right).
% "|" draws vertical rules and \hline draws horizontal rules.
% Convention: table captions go ABOVE the table. Always add a \label
% and refer to the table with \ref (e.g. Table~\ref{tab:setup}).
Table~\ref{tab:setup} summarises the configuration used in the experiments.

\begin{table}[h]
    \centering
    \caption{Example experimental configuration.}
    \label{tab:setup}
    \begin{tabular}{|l|c|r|}
        \hline
        \textbf{Parameter} & \textbf{Symbol} & \textbf{Value} \\ \hline
        Learning rate      & $\eta$          & 0.001          \\ \hline
        Batch size         & $B$             & 64             \\ \hline
        Epochs             & $N$             & 200            \\ \hline
    \end{tabular}
\end{table}

\lipsum[7]
